\typeout{NT FILE models-architecture.tex}%

\chapter{Models and architecture}
\label{cha:models-architecture}

In this chapter, we discuss the models and architectures used to classify
the fire risk of isolated houses. We explore the different models and
their respective architectures, as well as the classification layer used
to predict the fire risk. It will also be explained the unfreezing layer
strategy used when training the pre-trained models into our dataset.

\section{DenseNet-161 Architecture}

DenseNet-161 is a deep convolutional neural network that belongs to the DenseNet 
family, which was designed to optimize feature reuse and improve information flow 
between layers. The architecture follows the key principle of densely connected 
convolutional layers, where each layer receives input from all preceding layers, 
ensuring better gradient propagation and reducing redundancy in learned features.

\begin{figure}[ht]
    \centering
    \includegraphics[width=0.8\textwidth]{../Figures/DensenetArchitect.png}
    \caption{Representation of Densenet with 3 Dense blocks from \cite{DensenetResearchGate2025}.}
    \label{fig:densenet_block}
\end{figure}

One of the defining aspects of DenseNet-161 is the use of dense blocks, each 
consisting of multiple convolutional layers that share information through direct 
connections as we can see in Fig \ref{fig:densenet_block}. These dense connections mitigate the vanishing gradient problem and 
allow the model to achieve high accuracy with fewer parameters compared to 
traditional deep architectures. Between these dense blocks, transition layers 
are employed to compress feature maps, reducing computational cost while
maintaining effective feature extraction.

DenseNet-161 follows a structured layout that begins with an initial convolutional 
layer followed by four dense blocks. Each dense block consists of multiple layers 
applying $3 \times 3$ convolutions with batch normalization and ReLU activation. 
Within a dense block, each layer concatenates its feature maps with those of all 
preceding layers, increasing the model's representational capacity. Between the 
dense blocks, transition layers with $1 \times 1$ convolutions and average pooling 
are applied to control dimensionality.

The model takes an input of size $224 \times 224 \times 3$ and progressively reduces 
spatial dimensions while increasing depth. The final global average pooling layer 
extracts high-level features, which are then passed to a fully connected layer for 
classification.

\begin{table}[ht]
    \centering
    \begin{tabular}{|c|c|c|}
        \hline
        \textbf{Layer} & \textbf{Output Size} & \textbf{Details} \\ \hline
        Convolution & $112 \times 112 \times 96$ & $7 \times 7$, stride 2 \\ \hline
        Pooling & $56 \times 56 \times 96$ & $3 \times 3$ max pooling \\ \hline
        Dense Block 1 & $56 \times 56 \times 384$ & 6 layers \\ \hline
        Transition Layer 1 & $28 \times 28 \times 192$ & $1 \times 1$ conv + avg pool \\ \hline
        Dense Block 2 & $28 \times 28 \times 768$ & 12 layers \\ \hline
        Transition Layer 2 & $14 \times 14 \times 384$ & $1 \times 1$ conv + avg pool \\ \hline
        Dense Block 3 & $14 \times 14 \times 2112$ & 36 layers \\ \hline
        Transition Layer 3 & $7 \times 7 \times 1056$ & $1 \times 1$ conv + avg pool \\ \hline
        Dense Block 4 & $7 \times 7 \times 2208$ & 24 layers \\ \hline
        Pooling & $1 \times 1 \times 2208$ & Global average pooling \\ \hline
        Fully Connected & $1 \times 1 \times C$ & Classification layer \\ \hline
    \end{tabular}
    \caption{Layer-wise structure of DenseNet-161.}
    \label{tab:densenet_layers}
\end{table}

A key advantage of DenseNet-161 is its efficiency in parameter utilization. 
While traditional deep networks often suffer from redundancy, DenseNet leverages 
feature reuse to achieve higher accuracy with fewer parameters. This property 
makes it particularly effective in tasks requiring robust feature extraction 
with limited computational resources.

Additionally, the architecture incorporates residual-like connections without 
summation, preserving detailed information across layers. The presence of 
transition layers ensures that the number of feature maps does not grow 
excessively, keeping computational costs manageable.

DenseNet-161 has been widely adopted in image classification tasks due to its 
superior performance on benchmark datasets such as ImageNet. The combination 
of dense connectivity, batch normalization, and efficient pooling operations 
enables the model to learn highly discriminative features while maintaining 
computational efficiency. This makes it a strong candidate for our fire risk
classification task.

\section{ResNet-50 Architecture}

ResNet-50 is a deep convolutional neural network that belongs to the ResNet 
family, which introduced residual learning to significantly improve gradient 
propagation in deep architectures. The key innovation in ResNet models is 
the use of residual connections, which allow layers to learn residual mappings 
instead of direct transformations, mitigating the vanishing gradient problem 
and enabling the training of very deep networks.

The architecture of ResNet-50 consists of 50 layers, primarily composed 
of residual blocks. Each residual block includes identity shortcuts that 
bypass multiple convolutional layers, facilitating the flow of gradients 
through the network. This approach allows the model to learn deeper feature 
representations without suffering from degradation in accuracy due to increasing 
depth.

%\begin{figure}[ht]
%    \centering
%    \includegraphics[width=0.6\textwidth]{../Figures/ResNetBlock.png}
%    \caption{Representation of a residual block used in ResNet-50.}
%    \label{fig:resnet_block}
%\end{figure}

ResNet-50 follows a structured layout that begins with an initial convolutional 
layer followed by multiple residual blocks. Each residual block consists of three 
convolutional layers: a $1 \times 1$ convolution for dimensionality reduction, a 
$3 \times 3$ convolution for feature extraction, and another $1 \times 1$ convolution 
for dimensionality restoration. The shortcut connections add the original input 
back to the transformed output, improving feature reuse and allowing deeper 
networks to be trained efficiently.

The model takes an input of size $224 \times 224 \times 3$ and progressively 
reduces spatial dimensions while increasing feature depth. The final global 
average pooling layer extracts high-level features, which are then passed to 
a fully connected layer for classification.

\begin{table}[ht]
    \centering
    \begin{tabular}{|c|c|c|}
        \hline
        \textbf{Layer} & \textbf{Output Size} & \textbf{Details} \\ \hline
        Convolution & $112 \times 112 \times 64$ & $7 \times 7$, stride 2 \\ \hline
        Pooling & $56 \times 56 \times 64$ & $3 \times 3$ max pooling \\ \hline
        Residual Block 1 & $56 \times 56 \times 256$ & 3 layers $1\times1, 3\times3, 1\times1 \times 3$ \\ \hline
        Residual Block 2 & $28 \times 28 \times 512$ & 4 layers $1\times1, 3\times3, 1\times1 \times 4$ \\ \hline
        Residual Block 3 & $14 \times 14 \times 1024$ & 6 layers $1\times1, 3\times3, 1\times1 \times 6$ \\ \hline
        Residual Block 4 & $7 \times 7 \times 2048$ & 3 layers $1\times1, 3\times3, 1\times1 \times 3$ \\ \hline
        Pooling & $1 \times 1 \times 2048$ & Global average pooling \\ \hline
        Fully Connected & $1 \times 1 \times C$ & Classification layer \\ \hline
    \end{tabular}
    \caption{Layer-wise structure of ResNet-50.}
    \label{tab:resnet_layers}
\end{table}

A key advantage of ResNet-50 is its ability to train deep models without 
suffering from performance degradation. The residual connections enable efficient 
gradient flow, making the network more robust to vanishing gradients while 
significantly improving convergence speed. These properties make ResNet-50 
particularly effective for large-scale image classification tasks.

Additionally, ResNet-50 benefits from batch normalization and ReLU activation 
to enhance training stability and efficiency. Its hierarchical structure allows 
it to capture increasingly complex patterns as depth increases, making it 
well-suited for high-resolution image analysis.

%\begin{figure}[ht]
%    \centering
%    \includegraphics[width=0.6\textwidth]{../Figures/ResNetDiagram.png}
%    \caption{Overall architecture of ResNet-50, illustrating its residual connections and hierarchical feature extraction.}
%    \label{fig:resnet_overview}
%\end{figure}

ResNet-50 has demonstrated strong performance on benchmark datasets such as 
ImageNet and has been widely used in deep learning applications. Its combination 
of depth, residual learning, and efficient feature extraction makes it a powerful 
architecture for image classification, object detection, and other vision-related tasks.


\section{EfficientNet-B5 Architecture}

EfficientNet-B5 is a deep convolutional neural network that belongs to the EfficientNet family, which was designed to achieve high accuracy while maintaining computational efficiency. The EfficientNet models introduce a compound scaling method that balances depth, width, and resolution to improve performance without drastically increasing computational costs.

The key innovation behind EfficientNet is the use of the compound scaling principle, which scales the network's depth (number of layers), width (number of channels per layer), and resolution (input image size) proportionally, rather than arbitrarily increasing one of these factors. This results in a more optimized architecture that achieves higher accuracy while using fewer parameters compared to traditional deep networks.

%\begin{figure}[ht]
%    \centering
%    \includegraphics[width=0.6\textwidth]{../Figures/EfficientNetScaling.png}
%    \caption{Visualization of compound scaling used in EfficientNet models.}
%    \label{fig:efficientnet_scaling}
%\end{figure}

EfficientNet-B5 follows a structured design that consists of several Mobile Inverted 
Bottleneck Convolution (MBConv) blocks, originally introduced in MobileNetV2. 
These blocks use depthwise separable convolutions, reducing the number of 
parameters while maintaining feature extraction efficiency. Each MBConv block 
contains an expansion phase, a depthwise convolution, and a projection phase,
ensuring efficient computation and feature reuse.

The network begins with a standard convolutional layer followed by multiple 
MBConv blocks, gradually increasing feature depth while reducing spatial 
dimensions. The final layers consist of a global average pooling operation 
followed by a fully connected layer for classification.

\begin{table}[ht]
    \centering
    \begin{tabular}{|c|c|c|c|c|}
        \hline
        \textbf{Stage} & \textbf{Operator} & \textbf{Resolution} & \textbf{Channels} & \textbf{Repeats} \\ \hline
        Stem & Conv3x3 & $224 \times 224$ & 48 & 1 \\ \hline
        1 & MBConv1, k3x3 & $112 \times 112$ & 24 & 3 \\ \hline
        2 & MBConv6, k3x3 & $56 \times 56$ & 40 & 3 \\ \hline
        3 & MBConv6, k5x5 & $28 \times 28$ & 64 & 4 \\ \hline
        4 & MBConv6, k3x3 & $14 \times 14$ & 128 & 4 \\ \hline
        5 & MBConv6, k5x5 & $14 \times 14$ & 176 & 6 \\ \hline
        6 & MBConv6, k5x5 & $7 \times 7$ & 304 & 9 \\ \hline
        7 & MBConv6, k3x3 & $7 \times 7$ & 512 & 3 \\ \hline
        Head & Conv1x1 & $7 \times 7$ & 2048 & 1 \\ \hline
        Pooling & Global Average Pooling & $1 \times 1$ & 2048 & - \\ \hline
        Fully Connected & Dense Layer & $1 \times 1$ & C & - \\ \hline
    \end{tabular}
    \caption{Layer-wise structure of EfficientNet-B5.}
    \label{tab:efficientnet_layers}
\end{table}

One of the key advantages of EfficientNet-B5 is its ability to achieve high 
accuracy while maintaining computational efficiency. Compared to traditional 
deep networks, EfficientNet requires significantly fewer parameters while 
matching or surpassing their performance on benchmark datasets. This makes 
it an attractive choice for applications that require high precision without 
excessive computational overhead.

Additionally, EfficientNet-B5 incorporates Swish activation functions, batch 
normalization, and squeeze-and-excitation (SE) blocks, which adaptively 
recalibrate feature maps to improve learning efficiency. The network also 
supports higher resolution inputs ($456 \times 456$), enabling better feature 
extraction from detailed imagery.

%\begin{figure}[ht]
%    \centering
%    \includegraphics[width=0.6\textwidth]{../Figures/EfficientNetDiagram.png}
%    \caption{Overall architecture of EfficientNet-B5, illustrating its use of MBConv blocks and compound scaling.}
%    \label{fig:efficientnet_overview}
%\end{figure}

EfficientNet-B5 has demonstrated superior performance on ImageNet and other 
benchmark datasets, making it well-suited for applications requiring accurate 
and efficient image classification. Its combination of depthwise separable 
convolutions, squeeze-and-excitation blocks, and compound scaling allows it 
to achieve a strong balance between accuracy and computational cost.

\section{Computational Complexity and Model Size}
One of the key considerations when choosing a deep learning architecture is the balance between performance and computational efficiency. DenseNet-161, despite its deep architecture, maintains a relatively lower parameter count compared to other models like ResNet, owing to its feature reuse strategy.

Table \ref{tab:densenet_params} provides an overview of the model's computational footprint, including the number of parameters and approximate memory usage.

\begin{table}[ht]
    \centering
    \begin{tabular}{|c|c|c|}
        \hline
        \textbf{Model} & \textbf{Parameters} & \textbf{Size (FP32)} \\ \hline
        ResNet-50 & 25.6M & 98 MB \\ \hline
        DenseNet-161 & 28.7M & 115 MB \\ \hline
        EfficientNet-B5 & 30.0M & 120 MB \\ \hline
    \end{tabular}
    \caption{Comparison of model sizes and parameter counts.}
    \label{tab:densenet_params}
\end{table}



\section{Custom Classification Layer}

To adapt the pre-trained architectures for the fire risk classification task, a simple
custom classification layer was implemented, replacing the default fully connected 
(FC) layer of each model. The modified classification layer enhances the network's 
ability to generalize while providing additional flexibility in learning complex patterns.

The final layer of the pre-trained base models (DenseNet-161, ResNet-101, and 
EfficientNet-B5) was replaced with a fully connected classifier, structured as 
shown in Table \ref{tab:classification_layer}.

\begin{table}[ht]
    \centering
    \begin{tabular}{|c|c|c|}
        \hline
        \textbf{Layer} & \textbf{Size} & \textbf{Activation} \\ \hline
        Fully Connected (FC1) & Input: extracted features, Output: 4096 & ReLU \\ \hline
        Dropout & Probability: 0.4 & - \\ \hline
        Fully Connected (FC2) & Input: 4096, Output: 4096 & ReLU \\ \hline
        Dropout & Probability: 0.4 & - \\ \hline
        Output Layer (Binary Classification) & Input: 4096, Output: 1 & Sigmoid \\ \hline
        Output Layer (Multi-Class Classification) & Input: 4096, Output: Number of classes & Softmax \\ \hline
    \end{tabular}
    \caption{Structure of the custom classification layer.}
    \label{tab:classification_layer}
\end{table}

The final layer varies depending on the number of classes. For binary classification, 
a single output neuron with a Sigmoid activation function is used. For multi-class 
classification, a Softmax activation function is applied to produce class probabilities.

Different values of dropout rates were tested to assess their impact on regularization, 
as overfitting was a potential issue due to the relatively small dataset size. The chosen 
dropout probability of 0.4 was found to provide a balance between retaining relevant 
feature representations and preventing the model from memorizing the training data.

By introducing this custom classification layer, the network maintains the rich feature 
representations extracted by the deep feature extractor while enhancing its adaptability 
to the specific classification task.

\section{Ensemble Model}

\begin{table}[ht]
    \centering
    \begin{tabular}{|c|c|c|}
        \hline
        \textbf{Model} & \textbf{Parameters} & \textbf{Size (FP32)} \\ \hline
        ResNet-50 & 25.6M & 98 MB \\ \hline
        DenseNet-161 & 28.7M & 115 MB \\ \hline
        EfficientNet-B5 & 30.0M & 120 MB \\ \hline
    \end{tabular}
    \caption{Comparison of model sizes and parameter counts.}
    \label{tab:model_comparison2}
\end{table}
